% Options for packages loaded elsewhere
\PassOptionsToPackage{unicode}{hyperref}
\PassOptionsToPackage{hyphens}{url}
\documentclass[
  11pt,
]{article}
\usepackage{xcolor}
\usepackage[margin=1in]{geometry}
\usepackage{amsmath,amssymb}
\setcounter{secnumdepth}{-\maxdimen} % remove section numbering
\usepackage{iftex}
\ifPDFTeX
  \usepackage[T1]{fontenc}
  \usepackage[utf8]{inputenc}
  \usepackage{textcomp} % provide euro and other symbols
\else % if luatex or xetex
  \usepackage{unicode-math} % this also loads fontspec
  \defaultfontfeatures{Scale=MatchLowercase}
  \defaultfontfeatures[\rmfamily]{Ligatures=TeX,Scale=1}
\fi
\usepackage{lmodern}
\ifPDFTeX\else
  % xetex/luatex font selection
\fi
% Use upquote if available, for straight quotes in verbatim environments
\IfFileExists{upquote.sty}{\usepackage{upquote}}{}
\IfFileExists{microtype.sty}{% use microtype if available
  \usepackage[]{microtype}
  \UseMicrotypeSet[protrusion]{basicmath} % disable protrusion for tt fonts
}{}
\makeatletter
\@ifundefined{KOMAClassName}{% if non-KOMA class
  \IfFileExists{parskip.sty}{%
    \usepackage{parskip}
  }{% else
    \setlength{\parindent}{0pt}
    \setlength{\parskip}{6pt plus 2pt minus 1pt}}
}{% if KOMA class
  \KOMAoptions{parskip=half}}
\makeatother
\usepackage{color}
\usepackage{fancyvrb}
\newcommand{\VerbBar}{|}
\newcommand{\VERB}{\Verb[commandchars=\\\{\}]}
\DefineVerbatimEnvironment{Highlighting}{Verbatim}{commandchars=\\\{\}}
% Add ',fontsize=\small' for more characters per line
\newenvironment{Shaded}{}{}
\newcommand{\AlertTok}[1]{\textcolor[rgb]{1.00,0.00,0.00}{\textbf{#1}}}
\newcommand{\AnnotationTok}[1]{\textcolor[rgb]{0.38,0.63,0.69}{\textbf{\textit{#1}}}}
\newcommand{\AttributeTok}[1]{\textcolor[rgb]{0.49,0.56,0.16}{#1}}
\newcommand{\BaseNTok}[1]{\textcolor[rgb]{0.25,0.63,0.44}{#1}}
\newcommand{\BuiltInTok}[1]{\textcolor[rgb]{0.00,0.50,0.00}{#1}}
\newcommand{\CharTok}[1]{\textcolor[rgb]{0.25,0.44,0.63}{#1}}
\newcommand{\CommentTok}[1]{\textcolor[rgb]{0.38,0.63,0.69}{\textit{#1}}}
\newcommand{\CommentVarTok}[1]{\textcolor[rgb]{0.38,0.63,0.69}{\textbf{\textit{#1}}}}
\newcommand{\ConstantTok}[1]{\textcolor[rgb]{0.53,0.00,0.00}{#1}}
\newcommand{\ControlFlowTok}[1]{\textcolor[rgb]{0.00,0.44,0.13}{\textbf{#1}}}
\newcommand{\DataTypeTok}[1]{\textcolor[rgb]{0.56,0.13,0.00}{#1}}
\newcommand{\DecValTok}[1]{\textcolor[rgb]{0.25,0.63,0.44}{#1}}
\newcommand{\DocumentationTok}[1]{\textcolor[rgb]{0.73,0.13,0.13}{\textit{#1}}}
\newcommand{\ErrorTok}[1]{\textcolor[rgb]{1.00,0.00,0.00}{\textbf{#1}}}
\newcommand{\ExtensionTok}[1]{#1}
\newcommand{\FloatTok}[1]{\textcolor[rgb]{0.25,0.63,0.44}{#1}}
\newcommand{\FunctionTok}[1]{\textcolor[rgb]{0.02,0.16,0.49}{#1}}
\newcommand{\ImportTok}[1]{\textcolor[rgb]{0.00,0.50,0.00}{\textbf{#1}}}
\newcommand{\InformationTok}[1]{\textcolor[rgb]{0.38,0.63,0.69}{\textbf{\textit{#1}}}}
\newcommand{\KeywordTok}[1]{\textcolor[rgb]{0.00,0.44,0.13}{\textbf{#1}}}
\newcommand{\NormalTok}[1]{#1}
\newcommand{\OperatorTok}[1]{\textcolor[rgb]{0.40,0.40,0.40}{#1}}
\newcommand{\OtherTok}[1]{\textcolor[rgb]{0.00,0.44,0.13}{#1}}
\newcommand{\PreprocessorTok}[1]{\textcolor[rgb]{0.74,0.48,0.00}{#1}}
\newcommand{\RegionMarkerTok}[1]{#1}
\newcommand{\SpecialCharTok}[1]{\textcolor[rgb]{0.25,0.44,0.63}{#1}}
\newcommand{\SpecialStringTok}[1]{\textcolor[rgb]{0.73,0.40,0.53}{#1}}
\newcommand{\StringTok}[1]{\textcolor[rgb]{0.25,0.44,0.63}{#1}}
\newcommand{\VariableTok}[1]{\textcolor[rgb]{0.10,0.09,0.49}{#1}}
\newcommand{\VerbatimStringTok}[1]{\textcolor[rgb]{0.25,0.44,0.63}{#1}}
\newcommand{\WarningTok}[1]{\textcolor[rgb]{0.38,0.63,0.69}{\textbf{\textit{#1}}}}
\usepackage{longtable,booktabs,array}
\newcounter{none} % for unnumbered tables
\usepackage{calc} % for calculating minipage widths
% Correct order of tables after \paragraph or \subparagraph
\usepackage{etoolbox}
\makeatletter
\patchcmd\longtable{\par}{\if@noskipsec\mbox{}\fi\par}{}{}
\makeatother
% Allow footnotes in longtable head/foot
\IfFileExists{footnotehyper.sty}{\usepackage{footnotehyper}}{\usepackage{footnote}}
\makesavenoteenv{longtable}
\usepackage{graphicx}
\makeatletter
\newsavebox\pandoc@box
\newcommand*\pandocbounded[1]{% scales image to fit in text height/width
  \sbox\pandoc@box{#1}%
  \Gscale@div\@tempa{\textheight}{\dimexpr\ht\pandoc@box+\dp\pandoc@box\relax}%
  \Gscale@div\@tempb{\linewidth}{\wd\pandoc@box}%
  \ifdim\@tempb\p@<\@tempa\p@\let\@tempa\@tempb\fi% select the smaller of both
  \ifdim\@tempa\p@<\p@\scalebox{\@tempa}{\usebox\pandoc@box}%
  \else\usebox{\pandoc@box}%
  \fi%
}
% Set default figure placement to htbp
\def\fps@figure{htbp}
\makeatother
\setlength{\emergencystretch}{3em} % prevent overfull lines
\providecommand{\tightlist}{%
  \setlength{\itemsep}{0pt}\setlength{\parskip}{0pt}}
\usepackage{bookmark}
\IfFileExists{xurl.sty}{\usepackage{xurl}}{} % add URL line breaks if available
\urlstyle{same}
\hypersetup{
  hidelinks,
  pdfcreator={LaTeX via pandoc}}

\author{}
\date{}


\usepackage{seqsplit}
\tolerance=9999
\emergencystretch=3em
\hbadness=9999
\begin{document}

\section{Topological Mismatch in Distributed Inference: Mechanized
Models and Protocol-Level Evidence for Fork/Race/Fold
Scheduling}\label{topological-mismatch-in-distributed-inference-mechanized-models-and-protocol-level-evidence-for-forkracefold-scheduling}

\emph{\href{https://www.patreon.com/cw/twbuley}{Taylor William Buley} --
Independent Researcher}\\
\emph{\href{mailto:taylor@forkjoin.ai}{\nolinkurl{taylor@forkjoin.ai}}}

\subsection{Abstract}\label{abstract}

Path-like execution can impose avoidable coordination cost when the
workload's natural structure has multiple useful paths. This paper
studies that mismatch in a finite-DAG fork/race/fold model using the
diagnostic Delta\_beta = beta\_1* - beta\_1\^{}exec, computed from
bounded workload/execution graph pairs under a correctness envelope
C1-C4. The novelty claim is not pipelining, DAG scheduling, or stream
multiplexing by themselves. The main formal result is an adequacy
trichotomy for emitted Aeon Flux sites: nonpositive deficit iff a
lossless transport realization exists, zero deficit iff a tight lossless
realization exists, and positive deficit forces path collision and
information loss within the declared model class. The engineering
surface instantiates that model in distributed staged computation and
Aeon Flow transport, and a compiler-emitted WallingtonRotation witness
pair reports matched pathCount = 4, streamCount = 4, Delta\_beta = 0 and
collided pathCount = 4, streamCount = 2, Delta\_beta = 2 cases. In the
companion artifacts, 8/8 primary wall-clock cells pass with median
speedups 11.785x-21.620x, and an Aeon Flow versus HTTP/3 protocol corpus
over 144 sites and 12,371 resources shows a 72.252\% median framing gain
with positive completion-latency gains across all 6 primary
environments. The claim is local: bounded graph pairs and emitted
witnesses make Delta\_beta operational enough to distinguish structural
path collapse from within-class protocol cost.

\subsection{1. Problem and Claim}\label{problem-and-claim}

\textbf{Thesis.} Path-like execution topologies impose avoidable
coordination cost when the workload's natural structure has positive
parallel-path complexity.

The target is narrow: distributed inference and the transport surface
that carries it across machines. A staged inference workload can contain
branch points, joins, and independent paths that are correct to keep
alive simultaneously. A path-like runtime collapses that structure into
one ordered line. Sometimes that is the right topology. Sometimes it is
a mismatch that inserts waiting, framing waste, or unnecessary
synchronization even though the workload itself would allow more
structure.

This paper isolates that mismatch as a bounded structural question. The
diagnostic is \texttt{Delta\_beta}, the difference between workload-side
path demand and runtime-side path budget at one declared abstraction
layer. The main paper-side claim is stronger than ``parallelism helps.''
Within the declared Aeon Flux site model, \texttt{Delta\_beta} is an
obstruction variable:

\begin{enumerate}
\def\labelenumi{\arabic{enumi}.}
\tightlist
\item
  \texttt{Delta\_beta\ \textless{}=\ 0} iff a lossless transport
  realization exists.
\item
  \texttt{Delta\_beta\ =\ 0} iff a tight lossless transport realization
  exists.
\item
  \texttt{Delta\_beta\ \textgreater{}\ 0} forces path collision and
  information loss.
\end{enumerate}

That theorem family is the center of the flagship. The scheduler and
protocol sections exist to show that the obstruction is not only formal.
A staged runtime can expose useful paths and then lose them again at the
wire if the transport silently reserializes the flow. Conversely, a
stream-preserving transport cannot help if the scheduler never exposes
more than one useful path. The paper therefore studies one coupled
execution surface: scheduler plus transport.

The contribution is intentionally small.

\begin{enumerate}
\def\labelenumi{\arabic{enumi}.}
\tightlist
\item
  It gives a finite-DAG fork/race/fold model with an explicit
  correctness envelope C1-C4.
\item
  It operationalizes \texttt{Delta\_beta} on bounded workload/execution
  graph pairs instead of leaving it as prose.
\item
  It proves the adequacy trichotomy above for emitted Aeon Flux sites
  and shows a matched/collided emitted \texttt{WallingtonRotation}
  witness pair.
\item
  It tests the same bridge with bounded wall-clock evidence,
  protocol-corpus evidence, and a one-path boundary control.
\end{enumerate}

The paper therefore does not ask novelty credit for overlap, dataflow,
or stream multiplexing in isolation. If the adequacy trichotomy and the
emitted witness bridge were removed, the remaining artifact would be a
bounded scheduler-and-protocol note. The flagship claim, if it earns any
novelty credit at all, is narrower: the same deficit quantity appears as
a graph-pair diagnostic, the antecedent of an existence/impossibility
theorem, and the payload of a compiler-emitted witness consumed by the
paper.

The flagship stays short on purpose because the claim is local and
falsifiable. If the argument cannot survive as one theorem family, one
emitted witness, and two bounded gates, then it is not yet ready for
publication.

\subsection{2. Model and Notation}\label{model-and-notation}

\subsubsection{2.1 Finite DAG Scheduling
Surface}\label{finite-dag-scheduling-surface}

Let \texttt{G\ =\ (V,\ E)} be a finite DAG representing a bounded
execution slice. Nodes denote local computations or transport actions.
Edges denote causality. A fork node has out-degree greater than one, a
fold node has in-degree greater than one, and a linear chain has
in-degree and out-degree at most one. The model uses four primitives:

\begin{enumerate}
\def\labelenumi{\arabic{enumi}.}
\tightlist
\item
  \texttt{Fork(S,\ O)\ -\textgreater{}\ \{S\_1,\ ...,\ S\_k\}} creates
  branch-local states.
\item
  \texttt{Race(\{S\_i\})\ -\textgreater{}\ (S\_w,\ i\_w)} advances
  branches concurrently and accepts the first valid completion.
\item
  \texttt{Fold(\{S\_i\},\ f)\ -\textgreater{}\ S*} merges branch
  outcomes with a deterministic fold function \texttt{f}.
\item
  \texttt{Vent(S\_i)\ -\textgreater{}\ bottom} terminates a branch
  without corrupting its siblings.
\end{enumerate}

The paper does not need a larger primitive set. The local questions are
scheduling questions: when is a path split useful, when does a branch
need an all-branch fold, and when can an execution surface safely accept
early completion or local failure?

\subsubsection{2.2 Parallel-Path Complexity and
Delta\_beta}\label{parallel-path-complexity-and-delta_beta}

Let \texttt{beta\_1*} denote the workload-side path demand at the
abstraction layer under study, and let \texttt{beta\_1\^{}exec} denote
the path budget exposed by the execution topology at that same layer.
The diagnostic is

\[
\Delta_\beta = \beta_1^* - \beta_1^{\mathrm{exec}}.
\]

If \texttt{Delta\_beta\ =\ 0}, the schedule is structurally matched at
that layer. If \texttt{Delta\_beta\ \textgreater{}\ 0}, the execution
surface is narrower than the modeled workload. The diagnostic is
directional and local. It is not a universal system score.

To keep that claim falsifiable, the paper uses one explicit extraction
rule. Fix one abstraction layer and one bounded execution slice. Build a
workload graph \texttt{G*} whose nodes are bounded events and whose
edges are only the dependencies required for correctness at that layer.
Build an effective execution graph \texttt{G\^{}exec} for the same slice
after scheduler and transport ordering have been applied. Then compute

\[
\beta_1(G) = \max(0, |E| - |V| + \beta_0(G))
\]

on each graph, using the same event vocabulary on both sides. In this
paper \texttt{beta\_1} is not presented as a universal measure of
executable multiplicity in every model. It is the bounded graph-side
quantity used by the adequacy theorem and the emitted witness surface.

A worked example is enough to show the rule. A 10-stream fork/join
workload has \texttt{\textbar{}V\textbar{}\ =\ 12},
\texttt{\textbar{}E\textbar{}\ =\ 20}, and \texttt{beta\_0\ =\ 1}, so
\texttt{beta\_1*\ =\ 9}. If an ordered single-pipe transport forces the
same slice into a chain, then \texttt{\textbar{}E\textbar{}\ =\ 11},
\texttt{beta\_1\^{}exec\ =\ 0}, and \texttt{Delta\_beta\ =\ 9}. If the
transport preserves the same 10 branch-local streams through independent
recovery, then \texttt{beta\_1\^{}exec\ =\ 9} and
\texttt{Delta\_beta\ =\ 0}. This graph-pair rule is the operational
meaning of the diagnostic throughout the flagship.

The extraction rule is deliberately small.

\begin{enumerate}
\def\labelenumi{\arabic{enumi}.}
\tightlist
\item
  Declare one abstraction layer and one event vocabulary.
\item
  Build \texttt{G*} using only the dependencies required for correctness
  at that layer.
\item
  Build \texttt{G\^{}exec} for the same slice after scheduler and
  transport ordering have been applied.
\item
  Compute \texttt{beta\_1(G*)} and \texttt{beta\_1(G\^{}exec)} on the
  same bounded event set.
\item
  Report \texttt{Delta\_beta\ =\ beta\_1*\ -\ beta\_1\^{}exec}, and if
  the slice is an emitted Aeon Flux site, compare the resulting
  \texttt{pathCount} and \texttt{streamCount} against the adequacy
  regime.
\end{enumerate}

Three execution classes matter in this paper.

{\def\LTcaptype{none} % do not increment counter
\begin{longtable}[]{@{}
  >{\raggedright\arraybackslash}p{(\linewidth - 6\tabcolsep) * \real{0.2308}}
  >{\raggedright\arraybackslash}p{(\linewidth - 6\tabcolsep) * \real{0.2308}}
  >{\raggedleft\arraybackslash}p{(\linewidth - 6\tabcolsep) * \real{0.3077}}
  >{\raggedright\arraybackslash}p{(\linewidth - 6\tabcolsep) * \real{0.2308}}@{}}
\toprule\noalign{}
\begin{minipage}[b]{\linewidth}\raggedright
Execution class
\end{minipage} & \begin{minipage}[b]{\linewidth}\raggedright
Informal topology
\end{minipage} & \begin{minipage}[b]{\linewidth}\raggedleft
Typical \texttt{beta\_1\^{}exec} reading
\end{minipage} & \begin{minipage}[b]{\linewidth}\raggedright
Use in this paper
\end{minipage} \\
\midrule\noalign{}
\endhead
\bottomrule\noalign{}
\endlastfoot
Sequential path & one ordered line & \texttt{0} & boundary case \\
Chunked pipeline & stage-local concurrency without workload branching &
positive but bounded by stage occupancy & staged baseline \\
Fork/race/fold schedule & explicit branch, race, and deterministic
collapse surface & matched to positive \texttt{beta\_1*} when the
workload contains multiple useful paths & target topology \\
\end{longtable}
}

The flagship claim is not that every problem has positive
\texttt{beta\_1*}. Some workloads are truly path-like. For those,
\texttt{Delta\_beta\ =\ 0} under a sequential schedule and no theorem in
this paper predicts a structural win from adding concurrency.

The paper also does not ask \texttt{Delta\_beta} to replace
serial-fraction or occupancy language. Those metrics answer different
questions.

{\def\LTcaptype{none} % do not increment counter
\begin{longtable}[]{@{}
  >{\raggedright\arraybackslash}p{(\linewidth - 4\tabcolsep) * \real{0.3333}}
  >{\raggedright\arraybackslash}p{(\linewidth - 4\tabcolsep) * \real{0.3333}}
  >{\raggedright\arraybackslash}p{(\linewidth - 4\tabcolsep) * \real{0.3333}}@{}}
\toprule\noalign{}
\begin{minipage}[b]{\linewidth}\raggedright
Quantity
\end{minipage} & \begin{minipage}[b]{\linewidth}\raggedright
Useful for
\end{minipage} & \begin{minipage}[b]{\linewidth}\raggedright
What it does not localize
\end{minipage} \\
\midrule\noalign{}
\endhead
\bottomrule\noalign{}
\endlastfoot
serial fraction & classical scaling ceilings & whether the serial
bottleneck was demanded by the workload or inserted by the execution
surface \\
occupancy / utilization & idle-slot and saturation diagnosis & whether a
lower layer silently collapsed visible path structure \\
\texttt{Delta\_beta} & locating missing path budget on a bounded slice &
absolute system quality or a universal speedup scalar \\
\end{longtable}
}

\texttt{Delta\_beta} earns its keep when concurrency is visible in the
workload graph and absent in the effective execution graph at the same
time.

\subsubsection{2.3 C1-C4 Correctness
Envelope}\label{c1-c4-correctness-envelope}

The finite DAG model is used only inside the following correctness
envelope.

{\def\LTcaptype{none} % do not increment counter
\begin{longtable}[]{@{}
  >{\raggedright\arraybackslash}p{(\linewidth - 4\tabcolsep) * \real{0.3333}}
  >{\raggedright\arraybackslash}p{(\linewidth - 4\tabcolsep) * \real{0.3333}}
  >{\raggedright\arraybackslash}p{(\linewidth - 4\tabcolsep) * \real{0.3333}}@{}}
\toprule\noalign{}
\begin{minipage}[b]{\linewidth}\raggedright
Condition
\end{minipage} & \begin{minipage}[b]{\linewidth}\raggedright
Statement
\end{minipage} & \begin{minipage}[b]{\linewidth}\raggedright
Operational reading
\end{minipage} \\
\midrule\noalign{}
\endhead
\bottomrule\noalign{}
\endlastfoot
C1 & Constraint locality & Stage-local ordering is sufficient for global
correctness. \\
C2 & Branch isolation & A vented branch cannot corrupt sibling
branches. \\
C3 & Deterministic fold & The merge function is deterministic on the
admitted branch outputs. \\
C4 & Termination & Every branch completes, vents, or times out in finite
time. \\
\end{longtable}
}

These are not cosmetic assumptions. If C1 fails, the system may require
a stronger global order than the schedule preserves. If C2 fails, a
local rejection can leak damage sideways. If C3 fails, fold introduces
ambiguity rather than removing it. If C4 fails, the path budget never
collapses and the schedule cannot be treated as a bounded computation.

Within the formal companion, C1-C4 are checked in two ways. TLA+ models
make them explicit as state invariants or progress constraints on
bounded executions. Lean 4 theorem schemas supply the quantitative
consequences that depend on those conditions. The paper uses only the
part of that surface needed for the pipeline and protocol claims below.

\subsubsection{2.4 Pipeline Variables}\label{pipeline-variables}

For a prompt or batch with \texttt{P} items, chunk size \texttt{B},
stage count \texttt{N}, and chunk count

\[
Q = \lceil P / B \rceil,
\]

the chunked pipeline time under the modeled schedule is

\[
T_{\mathrm{pipe}} = Q + N - 1,
\]

while the fully serialized baseline is

\[
T_{\mathrm{seq}} = P \cdot N.
\]

The step-count speedup under the stated idealization is therefore

\[
\mathrm{speedup} = \frac{P \cdot N}{\lceil P/B \rceil + N - 1}.
\]

The model uses two explicit assumptions for this equation.

\begin{enumerate}
\def\labelenumi{\arabic{enumi}.}
\tightlist
\item
  A1: per-chunk stage service times are homogeneous across stages.
\item
  A2: inter-stage communication and synchronization cost is zero.
\end{enumerate}

Those assumptions matter. The step-count theorem is a model claim about
scheduling shape, not a promise that real deployments will achieve the
same scalar under arbitrary transport, compute heterogeneity, or loss.

\subsubsection{2.5 Reynolds Regimes}\label{reynolds-regimes}

The pipeline Reynolds number is defined as

\[
Re = N / Q.
\]

This quantity measures how dense the pipeline is in stages per chunk.
Low \texttt{Re} means many chunks relative to stage count, which pushes
the system toward steady-state occupancy. High \texttt{Re} means few
chunks relative to stage count, which magnifies ramp-up and drain
effects.

Figure 1 is a supporting model-side picture rather than a second
contribution claim. It carries the inverted-scaling workload curves and
the regime map induced by \texttt{Re}.

\begin{figure}
\centering
\includegraphics[width=\textwidth,keepaspectratio,alt={Figure 1. Inverted scaling and Reynolds regime map from the Chapter 17 companion artifacts. The same analytic surface supports the speedup ceiling, the speedup floor, and the regime thresholds used to classify laminar, transitional, and turbulent execution.}]{companion-tests/artifacts/ch17-inverted-scaling-reynolds-figure.png}
\caption{Figure 1. Inverted scaling and Reynolds regime map from the
Chapter 17 companion artifacts. The same analytic surface supports the
speedup ceiling, the speedup floor, and the regime thresholds used to
classify laminar, transitional, and turbulent execution.}
\end{figure}

\subsection{3. Formal Results}\label{formal-results}

\subsubsection{3.1 Correctness Under Finite DAG
Decomposition}\label{correctness-under-finite-dag-decomposition}

The first theorem surface is sufficiency, not universality. Under
finite-DAG decomposition and C1-C4, fork/race/fold scheduling can
preserve correctness without requiring a globally serialized path. The
paper does not claim a completeness theorem for all distributed systems.
It uses the narrower constructive claim that a bounded workload can be
decomposed into linear chains, fork points, and join points, and that
the scheduler may operate on that decomposition without changing the
result when the four conditions hold.

That surface answers the first question: is the richer schedule
well-defined? The remaining results address the second question: when
does the richer schedule matter?

\subsubsection{3.2 Pipeline Speedup Floor and
Sandwich}\label{pipeline-speedup-floor-and-sandwich}

The core scheduling theorem used here is
\texttt{THM-PIPELINE-SPEEDUP-FLOOR} from \texttt{Multiplexing.lean}.
Under A1-A2, the pipelined time never exceeds the serialized time:

\[
\lceil P/B \rceil + N - 1 \le P \cdot N.
\]

This is the lower side of the speedup sandwich. The upper side comes
from the large-\texttt{P} limit:

\[
\mathrm{speedup} \to B \cdot N \quad \text{as } P \to \infty.
\]

Together these yield the bounded interval

\[
1 \le \mathrm{speedup} \le B \cdot N.
\]

The value of the theorem is not that the ceiling is surprising. The
value is that the floor is explicit. A staged pipeline under the modeled
assumptions is structurally non-worse than the serialized baseline. If
the workload admits a pipeline and the deployment still measures a loss,
that loss is evidence against the assumptions or evidence that the
deployed schedule is not the one the model claims to analyze.

The companion packages both the floor and the strict-speedup refinement.
The strict version applies for \texttt{P\ \textgreater{}=\ 2} and
\texttt{N\ \textgreater{}=\ 2}, ruling out the degenerate single-item
and single-stage boundaries where the pipeline collapses to the path.

\subsubsection{3.3 Reynolds
Classification}\label{reynolds-classification}

The second theorem family is a supporting classifier for execution
regimes. The mechanized claim used here is operational:

\begin{enumerate}
\def\labelenumi{\arabic{enumi}.}
\tightlist
\item
  \texttt{Re\ \textless{}\ 1/3} is laminar.
\item
  \texttt{1/3\ \textless{}=\ Re\ \textless{}\ 2/3} is transitional.
\item
  \texttt{Re\ \textgreater{}=\ 2/3} is turbulent.
\end{enumerate}

The companion theorem names \texttt{reynolds\_regime\_exhaustive},
\texttt{classifyRegime}, and \texttt{turbulent\_collapse\_floor} are
used because they make the schedule-level consequence explicit: the
regime is not only a label, it constrains which fold policies are safe
and how much collapse work the runtime should expect to pay.

The practical use of the classification is modest. It turns the sentence
``this pipeline is too thin to keep stages busy'' into a falsifiable
condition on \texttt{N/Q}. The paper does not need a fluid identity
claim. It only needs the regime map as a schedule-selection tool for
interpreting Gate1.

\subsubsection{3.4 Aeon Flux Site
Adequacy}\label{aeon-flux-site-adequacy}

The central theorem family in the flagship is now the emitted-site
adequacy surface. For positive-stream-budget Aeon Flux sites, the
mechanized trichotomy is:

\begin{enumerate}
\def\labelenumi{\arabic{enumi}.}
\tightlist
\item
  \texttt{Delta\_beta\ \textless{}=\ 0} iff a lossless FlowFrame
  transport realization exists.
\item
  \texttt{Delta\_beta\ =\ 0} iff a tight lossless realization exists.
\item
  \texttt{Delta\_beta\ \textgreater{}\ 0} forces path collision and
  positive information loss for every realization in the model class.
\end{enumerate}

This is the point at which \texttt{Delta\_beta} stops being a
descriptive label and becomes an obstruction variable. Once
\texttt{pathCount} and \texttt{streamCount} are emitted from lowered GG,
the deficit is no longer being inferred from benchmark scalars or read
off diagrams by hand. It is a concrete input to a proved trichotomy.

The importance of that result is narrow but decisive. The paper no
longer asks the reader to believe that positive deficit merely ``sounds
like'' trouble. Within the declared site model, positive deficit means
lossless transport is impossible. Zero deficit means tight lossless
transport is possible. That is the exact coup de grace the flagship
needs.

That is also the boundary against the obvious relabeling objection.
Classical DAG scheduling already says that richer dependency structure
can matter. The claim here is not merely that the graph has branches. It
is that one emitted, bounded quantity is sufficient to state an iff
adequacy result for the declared site model and to classify concrete
lowered sources without hand interpretation.

\subsubsection{3.5 What the Formal Surface
Certifies}\label{what-the-formal-surface-certifies}

The formal surface certifies an existence/impossibility claim in the
declared site model. It does not certify wall-clock speedups, and it
does not certify Gate2 latency numbers. Those are measured. The theorem
family tells us which transport regimes are structurally possible, and
the gates tell us how the deployed scheduler and protocol behave inside
those regimes.

That distinction matters for the protocol story. In the flagship, the
transport-side formal anchor is the Aeon Flux site adequacy theorem plus
executable framing/stream tests. Gate2 then asks a narrower empirical
question: once stream preservation has been separated from gross path
collapse, does Aeon Flow still reduce residual framing and completion
cost relative to explicit \texttt{HTTP/3} baselines on the declared
corpus?

\subsection{4. Systems Instantiation}\label{systems-instantiation}

\subsubsection{4.1 Distributed Staged
Computation}\label{distributed-staged-computation}

The concrete scheduler studied in the companion is distributed staged
computation for inference prefill. Prompt chunks move across stage
nodes. Each node preserves stage-local order for the stream it owns,
while the overall runtime overlaps multiple chunks across stages. This
is exactly the sort of setting where C1 matters: correctness follows
from local ordering rather than from a single total order over all work
in flight.

At the implementation level the schedule is a chunked pipeline rather
than a speculative executor. A chunk enters stage 1, advances when that
stage is free, and continues through the remaining stages while later
chunks begin behind it. The modeled makespan is the
\texttt{Q\ +\ N\ -\ 1} surface from Section 2.4. The measured wall-clock
evidence later in the paper tests whether the deployed system preserves
enough of that shape to separate chunked execution from the serialized
baseline.

\subsubsection{4.2 Emitted Witness Pair from Lowered
GG}\label{emitted-witness-pair-from-lowered-gg}

The flagship now uses a matched pair of compiler-emitted witnesses
rather than only a verbal example. From two minimal
\texttt{WallingtonRotation} sources in the companion artifact directory,
the Gnosis/Aeon Forge exporter emits the following Aeon Flux site
witnesses:

{\def\LTcaptype{none} % do not increment counter
\begin{longtable}[]{@{}
  >{\raggedright\arraybackslash}p{(\linewidth - 12\tabcolsep) * \real{0.1250}}
  >{\raggedright\arraybackslash}p{(\linewidth - 12\tabcolsep) * \real{0.1250}}
  >{\raggedleft\arraybackslash}p{(\linewidth - 12\tabcolsep) * \real{0.1667}}
  >{\raggedleft\arraybackslash}p{(\linewidth - 12\tabcolsep) * \real{0.1667}}
  >{\raggedleft\arraybackslash}p{(\linewidth - 12\tabcolsep) * \real{0.1667}}
  >{\raggedright\arraybackslash}p{(\linewidth - 12\tabcolsep) * \real{0.1250}}
  >{\raggedright\arraybackslash}p{(\linewidth - 12\tabcolsep) * \real{0.1250}}@{}}
\toprule\noalign{}
\begin{minipage}[b]{\linewidth}\raggedright
Site
\end{minipage} & \begin{minipage}[b]{\linewidth}\raggedright
Primitive id
\end{minipage} & \begin{minipage}[b]{\linewidth}\raggedleft
\texttt{pathCount}
\end{minipage} & \begin{minipage}[b]{\linewidth}\raggedleft
\texttt{streamCount}
\end{minipage} & \begin{minipage}[b]{\linewidth}\raggedleft
\texttt{Delta\_beta}
\end{minipage} & \begin{minipage}[b]{\linewidth}\raggedright
Regime
\end{minipage} & \begin{minipage}[b]{\linewidth}\raggedright
Theorem ref
\end{minipage} \\
\midrule\noalign{}
\endhead
\bottomrule\noalign{}
\endlastfoot
matched site & \texttt{rotated\_output} & \texttt{4} & \texttt{4} &
\texttt{0} & \texttt{tight\_lossless\_transport\_exists} &
\texttt{aeon\_flux\_site\_zero\_deficit\_iff\_tight\_lossless\_transport} \\
collided site & \texttt{collapsed\_output} & \texttt{4} & \texttt{2} &
\texttt{2} & \texttt{collision\_and\_information\_loss\_forced} &
\texttt{aeon\_flux\_site\_positive\_deficit\_forces\_collision\_and\_information\_loss} \\
\end{longtable}
}

This table is intentionally plain. It is the concrete object the theorem
talks about. The paper no longer needs to ask the reader to imagine what
matched and collided sites look like. It points to lowered source files
and emitted witness packets for both.

\subsubsection{4.3 Transport Scope and Explicit
Baseline}\label{transport-scope-and-explicit-baseline}

The scheduler runs over Aeon Flow, a binary framing layer whose unit of
transport is a \texttt{FlowFrame} with a 10-byte header:

\begin{enumerate}
\def\labelenumi{\arabic{enumi}.}
\tightlist
\item
  \texttt{stream\_id}: 2 bytes
\item
  \texttt{sequence}: 4 bytes
\item
  \texttt{flags}: 1 byte
\item
  \texttt{length}: 3 bytes
\end{enumerate}

That fact matters because the protocol question here is not only
throughput. It is whether the transport preserves enough path structure
for the scheduler to remain visible at the wire. Gate2 uses an explicit
\texttt{HTTP/3} baseline, and the claim is narrow: once stream
preservation has already been separated from gross path collapse, Aeon
Flow may still reduce residual framing and completion cost.

\subsection{5. Evaluation}\label{evaluation}

\subsubsection{5.1 Evaluation Policy}\label{evaluation-policy}

Only artifact-backed quantitative claims are used in the flagship body.
The evaluation surface is intentionally small:

\begin{enumerate}
\def\labelenumi{\arabic{enumi}.}
\tightlist
\item
  one emitted witness pair from lowered GG;
\item
  bounded graph-pair and one-path controls;
\item
  Gate1 for wall-clock staged execution;
\item
  Gate2 for protocol-corpus framing and completion behavior.
\end{enumerate}

The paper only makes claims attached to an anchor, an observable, and a
failure condition.

{\def\LTcaptype{none} % do not increment counter
\begin{longtable}[]{@{}
  >{\raggedright\arraybackslash}p{(\linewidth - 6\tabcolsep) * \real{0.2500}}
  >{\raggedright\arraybackslash}p{(\linewidth - 6\tabcolsep) * \real{0.2500}}
  >{\raggedright\arraybackslash}p{(\linewidth - 6\tabcolsep) * \real{0.2500}}
  >{\raggedright\arraybackslash}p{(\linewidth - 6\tabcolsep) * \real{0.2500}}@{}}
\toprule\noalign{}
\begin{minipage}[b]{\linewidth}\raggedright
Paper-side claim
\end{minipage} & \begin{minipage}[b]{\linewidth}\raggedright
Anchor
\end{minipage} & \begin{minipage}[b]{\linewidth}\raggedright
Observable used here
\end{minipage} & \begin{minipage}[b]{\linewidth}\raggedright
What would count against the claim
\end{minipage} \\
\midrule\noalign{}
\endhead
\bottomrule\noalign{}
\endlastfoot
emitted sites can classify \texttt{Delta\_beta} by theorem regime rather
than by prose & Aeon Flux site adequacy theorems plus emitted witness
export & the \texttt{WallingtonRotation} witness packet & emitted
\texttt{pathCount}, \texttt{streamCount}, and theorem regime disagree \\
\texttt{Delta\_beta} can be computed on bounded graph pairs &
workload/execution graph comparison over explicit traces & exact
recovery of one-path, ordered-pipe-collapse, and per-stream-preserving
cases & graph pairs fail to separate \texttt{Delta\_beta\ =\ 0} from
positive-deficit cases \\
chunked staged execution can remove path-imposed coordination cost under
A1-A2 & pipeline floor and sandwich theorem family & Gate1 latency
separation between serialized and chunked execution & chunked execution
fails to separate from, or reverses against, the serialized baseline on
the primary cells \\
once stream preservation is already present, protocol structure still
changes residual coordination cost & FlowFrame protocol checks plus
Gate2 corpus harness (\texttt{Aeon\ Flow} versus \texttt{HTTP/3}) &
Gate2 framing and completion behavior across the primary environment
matrix & framing or completion gains collapse against the explicit
HTTP/3 baseline \\
\end{longtable}
}

\subsubsection{5.2 Emitted Witness and Boundary
Checks}\label{emitted-witness-and-boundary-checks}

Before reading the deployment gates, the paper fixes the diagnostic on
explicit witnesses and graph pairs.

{\def\LTcaptype{none} % do not increment counter
\begin{longtable}[]{@{}
  >{\raggedright\arraybackslash}p{(\linewidth - 8\tabcolsep) * \real{0.1667}}
  >{\raggedleft\arraybackslash}p{(\linewidth - 8\tabcolsep) * \real{0.2222}}
  >{\raggedleft\arraybackslash}p{(\linewidth - 8\tabcolsep) * \real{0.2222}}
  >{\raggedleft\arraybackslash}p{(\linewidth - 8\tabcolsep) * \real{0.2222}}
  >{\raggedright\arraybackslash}p{(\linewidth - 8\tabcolsep) * \real{0.1667}}@{}}
\toprule\noalign{}
\begin{minipage}[b]{\linewidth}\raggedright
Case
\end{minipage} & \begin{minipage}[b]{\linewidth}\raggedleft
\texttt{beta\_1*} or \texttt{pathCount}
\end{minipage} & \begin{minipage}[b]{\linewidth}\raggedleft
\texttt{beta\_1\^{}exec} or \texttt{streamCount}
\end{minipage} & \begin{minipage}[b]{\linewidth}\raggedleft
\texttt{Delta\_beta}
\end{minipage} & \begin{minipage}[b]{\linewidth}\raggedright
Reading
\end{minipage} \\
\midrule\noalign{}
\endhead
\bottomrule\noalign{}
\endlastfoot
emitted matched site & \texttt{4} & \texttt{4} & \texttt{0} & tight
zero-deficit witness from lowered GG \\
emitted collided site & \texttt{4} & \texttt{2} & \texttt{2} &
positive-deficit witness; the emitted site itself falls into the
collision regime \\
one-path chain & \texttt{0} & \texttt{0} & \texttt{0} & exact parity
boundary; no structural win is predicted \\
10-stream fork/join over one ordered pipe & \texttt{9} & \texttt{0} &
\texttt{9} & positive mismatch; the transport has collapsed visible
concurrency back into one path \\
10-stream fork/join over per-stream recovery (\texttt{HTTP/3} or
\texttt{Aeon\ Flow}) & \texttt{9} & \texttt{9} & \texttt{0} & stream
independence is preserved; remaining differences are protocol cost \\
\end{longtable}
}

This is the direct operational check for the diagnostic.
\texttt{Delta\_beta} is not being reverse-engineered from Gate1 or Gate2
outcomes. It is being computed on bounded graph pairs first and emitted
witness packets second. The matched/collided witness pair matters
because it shows the adequacy theorem classifying concrete lowered
sources on both sides of the boundary instead of only on the
zero-deficit side.

The same surface also predicts the no-win boundary. In the companion,
the one-path control is realized as single-chunk execution
\texttt{Q\ =\ 1}, and the pipeline equation recovers the serial
\texttt{N}-step path:

\[
T_{\mathrm{pipe}} = Q + N - 1 = 1 + N - 1 = N.
\]

There is no hidden speedup term left once the path budget has collapsed
to one. This is a control, not a headline result. It shows that the
flagship is allowed to predict parity.

\subsubsection{5.3 Gate1: External Multi-Host Wall-Clock
Evidence}\label{gate1-external-multi-host-wall-clock-evidence}

Gate1 measures chunked versus serialized execution on an external
multi-host setup. The flagship uses the external multi-host report
because it is the cleanest bridge from the formal pipeline model to
deployed staged execution. The artifact uses two workloads, multiple RTT
and loss settings, bootstrap confidence intervals, and six distinct
endpoint hosts.

The result is simple and strong. 8/8 primary Gate1 cells pass. Across
those primary cells, the median speedup range is 11.785x-21.620x. The
minimum primary-cell speedup CI low is 11.365x. Across the full ten-cell
matrix, the minimum latency-improvement CI low is 3561 ms.

Those numbers matter for two reasons. First, they are comfortably above
the speedup floor. The deployment is not merely avoiding regressions; it
is showing a large, repeated separation from the serialized baseline.
Second, the range widens as the workload gets larger in stage count and
token count, which is exactly the direction predicted by the modeled
step-count surface.

The Gate1 figure is the most direct empirical view of the scheduling
claim.

Figure 2 shows the deployed wall-clock separation that corresponds to
the model-side speedup surface.

\begin{figure}
\centering
\includegraphics[width=\textwidth,keepaspectratio,alt={Figure 2. Gate1 wall-clock evidence from the Chapter 17 companion artifacts. The chart summarizes the ten-cell external multi-host matrix and shows the separation between serialized and chunked execution together with bootstrap speedup intervals.}]{companion-tests/artifacts/ch17-gate1-wallclock-figure.png}
\caption{Figure 2. Gate1 wall-clock evidence from the Chapter 17
companion artifacts. The chart summarizes the ten-cell external
multi-host matrix and shows the separation between serialized and
chunked execution together with bootstrap speedup intervals.}
\end{figure}

The flagship does not claim that the Gate1 scalar is a universal
constant. The point is narrower: on a bounded external deployment
surface, the measured behavior matches the qualitative prediction of the
formal scheduler. A path-like baseline is paying coordination cost that
the staged schedule avoids.

\subsubsection{5.4 Gate2: Protocol-Corpus
Evidence}\label{gate2-protocol-corpus-evidence}

Gate2 tests the transport side of the bridge on a seeded heterogeneous
corpus. The corpus contains 144 sites and 12,371 resources, and the six
primary environments span increasing RTT, decreasing bandwidth, and
rising loss. The explicit baseline is \texttt{HTTP/3}. That distinction
matters. Gate2 is not where the paper establishes the positive-deficit
ordered-pipe collapse case; the companion boundary checks already do
that on bounded graph pairs. Gate2 asks the narrower question of
within-class protocol efficiency: once the transport is already
stream-preserving, do Aeon Flow's framing and fold semantics still
improve framing and completion behavior relative to \texttt{HTTP/3}
across the declared impairment matrix?

Again the result is compact. 6/6 primary Gate2 cells pass. The framing
median gain is 72.252\%, and the minimum primary-cell framing CI low is
72.19\%. At the benign primary edge, the \texttt{rtt12-bw80-loss0}
environment still yields 23.62 ms median completion gain and 23.79 ms
p95 completion gain. At the harsh primary edge,
\texttt{rtt110-bw7-loss2pct} yields 92.39 ms median completion gain and
110.5 ms p95 completion gain.

The framing gain is nearly invariant across environments because it is
dominated by the declared framing models rather than by network
impairment. The completion gains are the network-sensitive part, and
they widen as impairment increases. The flagship therefore reads Gate2
narrowly and honestly: after the path-collapse boundary has been
separated elsewhere, protocol structure still changes the residual
coordination cost of carrying the staged workload over the wire.

Figure 3 makes the split visible: nearly constant framing advantage,
widening completion advantage.

\begin{figure}
\centering
\includegraphics[width=\textwidth,keepaspectratio,alt={Figure 3. Gate2 protocol-corpus evidence from the Chapter 17 companion artifacts. The chart separates framing gains from completion gains across the heterogeneous site corpus and shows that the protocol advantage persists across all six primary environments.}]{companion-tests/artifacts/ch17-gate2-protocol-corpus-figure.png}
\caption{Figure 3. Gate2 protocol-corpus evidence from the Chapter 17
companion artifacts. The chart separates framing gains from completion
gains across the heterogeneous site corpus and shows that the protocol
advantage persists across all six primary environments.}
\end{figure}

\subsubsection{5.5 Interpreting the Evaluation Surface
Together}\label{interpreting-the-evaluation-surface-together}

The emitted witness, graph-pair checks, Gate1, and Gate2 do different
jobs.

\begin{enumerate}
\def\labelenumi{\arabic{enumi}.}
\tightlist
\item
  The emitted witness shows that a lowered site can be assigned a
  theorem regime directly.
\item
  The graph-pair checks separate one-path, ordered-pipe-collapse, and
  per-stream-preserving cases before any deployment scalar is discussed.
\item
  Gate1 checks the staged scheduler against the serialized baseline.
\item
  Gate2 checks Aeon Flow against the explicit \texttt{HTTP/3} baseline
  once the gross collapse-vs-preservation question has already been
  isolated.
\item
  The one-path control checks that the same surface predicts the no-gain
  case when the path budget really is one.
\end{enumerate}

Taken together they support one bounded conclusion: \texttt{Delta\_beta}
is operational enough to separate positive-deficit path-collapse cases
from matched cases, and on that foundation the staged runtime plus
transport surface shows artifact-backed reductions in observed
coordination cost. That is the precise role of \texttt{Delta\_beta} in
the flagship: not a one-number oracle, but a way to state where missing
path budget is being inserted.

\subsection{6. Limitations and
Non-Claims}\label{limitations-and-non-claims}

The flagship deliberately narrows scope. The following boundaries are
part of the claim, not footnotes to it.

\begin{enumerate}
\def\labelenumi{\arabic{enumi}.}
\tightlist
\item
  The model is a finite DAG model. Infinite-state, adversarially
  adaptive, or unrestricted asynchronous systems are outside the direct
  theorem surface used here.
\item
  The graph-side \texttt{beta\_1} used here is a bounded site metric,
  not a claim that one topological object canonically measures
  executable multiplicity in every system model.
\item
  The emitted witness path shown in the flagship is currently scoped to
  lowered \texttt{WallingtonRotation} sites. The theorem family is
  broader than that one example, but the paper only needs one emitted
  witness.
\item
  The pipeline speedup formula is a modeled step-count claim under
  A1-A2. It is not a universal wall-clock predictor.
\item
  \texttt{Delta\_beta} is an obstruction variable inside the declared
  model, not a one-number oracle for system quality.
\item
  Gate1 is a bounded external multi-host matrix, not a claim about every
  deployment.
\item
  Gate2 is a seeded heterogeneous corpus with explicit impairment
  environments, and its protocol comparison is Aeon Flow versus
  \texttt{HTTP/3} rather than a universal transport shootout. Its
  framing gain is partly structural from declared framing models; its
  completion gains are the network-sensitive measurements.
\item
  The paper does not claim universal superiority of richer schedules.
  For small workloads, extremely fast serial paths, or high coordination
  overhead, the correct topology may be \texttt{beta\_1\ =\ 0}.
\item
  No body claim in this paper depends on extra-domain correspondence
  surfaces outside distributed systems and bounded protocol models.
\end{enumerate}

These limitations sharpen the result rather than weaken it. The claim is
strong precisely because it stays local.

\subsection{7. Related Work}\label{related-work}

This section is intentionally not a survey. It covers only the
literatures needed to identify the incremental claim of the flagship.

The speedup formulas in this paper sit beside, not above, the classical
scaling laws. Amdahl studies fixed-workload speed limits under an
explicit serial fraction, while Gustafson reframes the bound under
scaled workloads {[}1, 2{]}. Those laws remain the right tools for
ceiling questions. The current paper adds a different question: when is
part of the apparent serial bottleneck being inserted by the execution
surface rather than demanded by the workload?

The scheduling model is likewise adjacent to older work on explicit
concurrency surfaces. Kahn process networks, early dataflow
architectures, micropipelines, and MapReduce all make it natural to
reason over bounded dependency graphs and fork/fold structure rather
than one total order {[}3-6{]}. The flagship does not claim novelty for
DAG execution itself. Its narrower claim is that separating
workload-side path demand from runtime-side path budget, then binding
that gap to an adequacy theorem and an emitted witness, makes certain
failures easier to localize.

Put differently, prior work already supplies the language of explicit
concurrency. The flagship only claims one extra step beyond that
language: the deficit is not merely descriptive but theorem-bearing in
the emitted site model. If that step is uninteresting, then the paper
should be read as a small notation result rather than as a new scheduler
theory.

The protocol side belongs to a familiar lineage as well. HTTP/2 reduces
framing waste and introduces application-layer multiplexing, while QUIC
moves recovery and ordering decisions closer to individual streams {[}8,
9{]}. The paper does not claim to supersede that literature. Its
narrower claim is that, for the scheduler-plus-transport bridge studied
here, it is useful to distinguish gross path preservation from
within-class protocol cost instead of blending both questions into one
benchmark.

Finally, the formal-methods layer is deliberately conventional in the
best sense. TLA+ and related specification practice already provide a
language for bounded safety and progress arguments {[}7{]}. The novelty
claim here is not ``formal methods for systems'' in general. It is the
tighter connection between one bounded theorem family, one emitted
witness surface, and a small set of measured observables, all backed by
a public rerun surface rather than by prose alone {[}10{]}.

\subsection{8. Conclusion}\label{conclusion}

Path-like execution topologies impose avoidable coordination cost when
the workload's natural structure has positive parallel-path complexity.

That sentence is the whole paper. The finite-DAG model provides the
vocabulary. C1-C4 provide the correctness envelope. The pipeline floor
and Reynolds classification provide the scheduler-side quantitative
surface. The Aeon Flux site adequacy theorem turns \texttt{Delta\_beta}
into an obstruction result. The emitted \texttt{WallingtonRotation}
witness shows the theorem on a concrete lowered source. Gate1 and Gate2
then show that the bridge is not only formal: the staged runtime and the
protocol layer both exhibit artifact-backed gains on bounded evaluation
surfaces.

The result is narrow, and that is the point. One does not need a
universal theory of computation to say something useful about
topological mismatch in distributed inference. One needs one theorem
family, one emitted witness, explicit assumptions, and measured
evidence. If the theorem-plus-witness bridge is judged unimportant, the
paper fails on novelty. If it is judged useful, the claim is already as
small as it should be.

\subsection{Appendix A: Theorem-to-Artifact
Map}\label{appendix-a-theorem-to-artifact-map}

{\def\LTcaptype{none} % do not increment counter
\begin{longtable}[]{@{}
  >{\raggedright\arraybackslash}p{(\linewidth - 4\tabcolsep) * \real{0.3333}}
  >{\raggedright\arraybackslash}p{(\linewidth - 4\tabcolsep) * \real{0.3333}}
  >{\raggedright\arraybackslash}p{(\linewidth - 4\tabcolsep) * \real{0.3333}}@{}}
\toprule\noalign{}
\begin{minipage}[b]{\linewidth}\raggedright
Local claim in flagship
\end{minipage} & \begin{minipage}[b]{\linewidth}\raggedright
Theorem or artifact family
\end{minipage} & \begin{minipage}[b]{\linewidth}\raggedright
Concrete artifact
\end{minipage} \\
\midrule\noalign{}
\endhead
\bottomrule\noalign{}
\endlastfoot
finite DAG scheduling under explicit decomposition assumptions &
executable finite-DAG decomposition surface &
\texttt{companion-\allowbreak{}tests/\allowbreak{}src/\allowbreak{}dag-\allowbreak{}completeness.\allowbreak{}test.\allowbreak{}ts} \\
C1-C4 correctness envelope & bounded TLA+ and Lean scheduler surface &
\texttt{TemporalModel} companion modules plus \texttt{@a0n/aeon-logic}
rerun path \\
pipeline floor \texttt{T\_pipe\ \textless{}=\ T\_seq} &
\texttt{THM-PIPELINE-SPEEDUP-FLOOR} &
\texttt{formal/lean/Multiplexing.lean} \\
speedup interval \texttt{1\ \textless{}=\ speedup\ \textless{}=\ B*N} &
\texttt{pipeline\_speedup\_floor}, \texttt{pipeline\_strict\_speedup},
\texttt{pipeline\_speedup\_sandwich} &
\texttt{formal/lean/Multiplexing.lean},
\texttt{formal/lean/PredictionsRound15.lean} \\
Reynolds regime cut points & \texttt{reynolds\_regime\_exhaustive},
\texttt{classifyRegime}, \texttt{turbulent\_collapse\_floor} &
\texttt{formal/lean/ReynoldsBFT.lean},
\texttt{companion-\allowbreak{}tests/\allowbreak{}artifacts/\allowbreak{}ch17-\allowbreak{}inverted-\allowbreak{}scaling-\allowbreak{}reynolds-\allowbreak{}figure.\allowbreak{}\{json,md,png,svg\}} \\
Aeon Flux site adequacy trichotomy &
\texttt{aeon\_flux\_site\_nonpositive\_deficit\_iff\_lossless\_transport},
\texttt{aeon\_flux\_site\_zero\_deficit\_iff\_tight\_lossless\_transport},
\texttt{aeon\_flux\_site\_positive\_deficit\_forces\_collision\_and\_information\_loss}
&
\texttt{formal/\allowbreak{}lean/\allowbreak{}Lean/\allowbreak{}ForkRaceFoldTheorems/\allowbreak{}AeonFluxSiteAdequacy.\allowbreak{}lean},
\texttt{open-\allowbreak{}source/\allowbreak{}gnosis/\allowbreak{}examples/\allowbreak{}proofs/\allowbreak{}AeonFluxSiteAdequacy-\allowbreak{}lean.\allowbreak{}gg} \\
emitted matched/collided witness pair from lowered GG & Gnosis/Aeon
Forge site witness export &
\texttt{companion-\allowbreak{}tests/\allowbreak{}artifacts/\allowbreak{}ch17-\allowbreak{}wallington-\allowbreak{}rotation-\allowbreak{}site.\allowbreak{}gg},
\texttt{companion-\allowbreak{}tests/\allowbreak{}artifacts/\allowbreak{}ch17-\allowbreak{}wallington-\allowbreak{}rotation-\allowbreak{}site-\allowbreak{}witness.\allowbreak{}\{json,md\}},
\texttt{companion-\allowbreak{}tests/\allowbreak{}artifacts/\allowbreak{}ch17-\allowbreak{}wallington-\allowbreak{}rotation-\allowbreak{}positive-\allowbreak{}deficit-\allowbreak{}site.\allowbreak{}gg},
\texttt{companion-\allowbreak{}tests/\allowbreak{}artifacts/\allowbreak{}ch17-\allowbreak{}wallington-\allowbreak{}rotation-\allowbreak{}positive-\allowbreak{}deficit-\allowbreak{}site-\allowbreak{}witness.\allowbreak{}\{json,md\}} \\
operational \texttt{Delta\_beta} extraction from bounded graph pairs &
workload/execution graph comparison over explicit traces &
\texttt{companion-\allowbreak{}tests/\allowbreak{}src/\allowbreak{}deficit-\allowbreak{}evidence.\allowbreak{}test.\allowbreak{}ts} \\
10-byte FlowFrame header and stream semantics & framing law plus runtime
protocol tests & \texttt{companion-\allowbreak{}tests/\allowbreak{}src/\allowbreak{}flow-\allowbreak{}protocol.\allowbreak{}test.\allowbreak{}ts} \\
one-path boundary recovers sequential execution & explicit
\texttt{beta\_1\ =\ 0} boundary check &
\texttt{companion-\allowbreak{}tests/\allowbreak{}src/\allowbreak{}pipeline-\allowbreak{}topology.\allowbreak{}test.\allowbreak{}ts} \\
deployed scheduler evidence & Gate1 wall-clock matrix &
\texttt{companion-\allowbreak{}tests/\allowbreak{}artifacts/\allowbreak{}gate1-\allowbreak{}wallclock-\allowbreak{}external-\allowbreak{}multihost.\allowbreak{}json},
\texttt{companion-\allowbreak{}tests/\allowbreak{}artifacts/\allowbreak{}ch17-\allowbreak{}gate1-\allowbreak{}wallclock-\allowbreak{}figure.\allowbreak{}\{json,md,png,svg\}} \\
deployed protocol evidence & Gate2 protocol corpus &
\texttt{companion-\allowbreak{}tests/\allowbreak{}artifacts/\allowbreak{}gate2-\allowbreak{}protocol-\allowbreak{}corpus.\allowbreak{}json},
\texttt{companion-\allowbreak{}tests/\allowbreak{}artifacts/\allowbreak{}ch17-\allowbreak{}gate2-\allowbreak{}protocol-\allowbreak{}corpus-\allowbreak{}figure.\allowbreak{}\{json,md,png,svg\}} \\
\end{longtable}
}

\subsection{Appendix B: Reproduction
Surface}\label{appendix-b-reproduction-surface}

The flagship paper has its own build and validation path so that it can
be checked independently from the full theory chapter.

\subsubsection{Build}\label{build}

\begin{Shaded}
\begin{Highlighting}[]
\ExtensionTok{pnpm}\NormalTok{ run manuscript:arxiv:flagship:tex}
\ExtensionTok{pnpm}\NormalTok{ run manuscript:arxiv:flagship:build}
\ExtensionTok{pnpm}\NormalTok{ run manuscript:arxiv:flagship:package}
\end{Highlighting}
\end{Shaded}

\subsubsection{Emitted witness rerun}\label{emitted-witness-rerun}

\begin{Shaded}
\begin{Highlighting}[]
\ExtensionTok{pnpm} \AttributeTok{{-}{-}dir}\NormalTok{ apps/aeon{-}forge exec bun src/cli/cli.tsx topology analyze }\DataTypeTok{\textbackslash{}}
\NormalTok{  /Users/buley/Documents/Code/emotions/open{-}source/aeon/docs/ebooks/145{-}log{-}rolling{-}pipelined{-}prefill/companion{-}tests/artifacts/ch17{-}wallington{-}rotation{-}site.gg }\DataTypeTok{\textbackslash{}}
  \AttributeTok{{-}{-}format}\NormalTok{ summary}

\ExtensionTok{pnpm} \AttributeTok{{-}{-}dir}\NormalTok{ apps/aeon{-}forge exec bun src/cli/cli.tsx topology export }\DataTypeTok{\textbackslash{}}
\NormalTok{  /Users/buley/Documents/Code/emotions/open{-}source/aeon/docs/ebooks/145{-}log{-}rolling{-}pipelined{-}prefill/companion{-}tests/artifacts/ch17{-}wallington{-}rotation{-}site.gg }\DataTypeTok{\textbackslash{}}
  \AttributeTok{{-}{-}format}\NormalTok{ json }\DataTypeTok{\textbackslash{}}
  \AttributeTok{{-}{-}out}\NormalTok{ /Users/buley/Documents/Code/emotions/open{-}source/aeon/docs/ebooks/145{-}log{-}rolling{-}pipelined{-}prefill/companion{-}tests/artifacts/ch17{-}wallington{-}rotation{-}site{-}witness.json}

\ExtensionTok{pnpm} \AttributeTok{{-}{-}dir}\NormalTok{ apps/aeon{-}forge exec bun src/cli/cli.tsx topology export }\DataTypeTok{\textbackslash{}}
\NormalTok{  /Users/buley/Documents/Code/emotions/open{-}source/aeon/docs/ebooks/145{-}log{-}rolling{-}pipelined{-}prefill/companion{-}tests/artifacts/ch17{-}wallington{-}rotation{-}site.gg }\DataTypeTok{\textbackslash{}}
  \AttributeTok{{-}{-}format}\NormalTok{ markdown }\DataTypeTok{\textbackslash{}}
  \AttributeTok{{-}{-}out}\NormalTok{ /Users/buley/Documents/Code/emotions/open{-}source/aeon/docs/ebooks/145{-}log{-}rolling{-}pipelined{-}prefill/companion{-}tests/artifacts/ch17{-}wallington{-}rotation{-}site{-}witness.md}

\ExtensionTok{pnpm} \AttributeTok{{-}{-}dir}\NormalTok{ apps/aeon{-}forge exec bun src/cli/cli.tsx topology analyze }\DataTypeTok{\textbackslash{}}
\NormalTok{  /Users/buley/Documents/Code/emotions/open{-}source/aeon/docs/ebooks/145{-}log{-}rolling{-}pipelined{-}prefill/companion{-}tests/artifacts/ch17{-}wallington{-}rotation{-}positive{-}deficit{-}site.gg }\DataTypeTok{\textbackslash{}}
  \AttributeTok{{-}{-}format}\NormalTok{ summary}

\ExtensionTok{pnpm} \AttributeTok{{-}{-}dir}\NormalTok{ apps/aeon{-}forge exec bun src/cli/cli.tsx topology export }\DataTypeTok{\textbackslash{}}
\NormalTok{  /Users/buley/Documents/Code/emotions/open{-}source/aeon/docs/ebooks/145{-}log{-}rolling{-}pipelined{-}prefill/companion{-}tests/artifacts/ch17{-}wallington{-}rotation{-}positive{-}deficit{-}site.gg }\DataTypeTok{\textbackslash{}}
  \AttributeTok{{-}{-}format}\NormalTok{ json }\DataTypeTok{\textbackslash{}}
  \AttributeTok{{-}{-}out}\NormalTok{ /Users/buley/Documents/Code/emotions/open{-}source/aeon/docs/ebooks/145{-}log{-}rolling{-}pipelined{-}prefill/companion{-}tests/artifacts/ch17{-}wallington{-}rotation{-}positive{-}deficit{-}site{-}witness.json}

\ExtensionTok{pnpm} \AttributeTok{{-}{-}dir}\NormalTok{ apps/aeon{-}forge exec bun src/cli/cli.tsx topology export }\DataTypeTok{\textbackslash{}}
\NormalTok{  /Users/buley/Documents/Code/emotions/open{-}source/aeon/docs/ebooks/145{-}log{-}rolling{-}pipelined{-}prefill/companion{-}tests/artifacts/ch17{-}wallington{-}rotation{-}positive{-}deficit{-}site.gg }\DataTypeTok{\textbackslash{}}
  \AttributeTok{{-}{-}format}\NormalTok{ markdown }\DataTypeTok{\textbackslash{}}
  \AttributeTok{{-}{-}out}\NormalTok{ /Users/buley/Documents/Code/emotions/open{-}source/aeon/docs/ebooks/145{-}log{-}rolling{-}pipelined{-}prefill/companion{-}tests/artifacts/ch17{-}wallington{-}rotation{-}positive{-}deficit{-}site{-}witness.md}
\end{Highlighting}
\end{Shaded}

\subsubsection{Targeted manuscript
validation}\label{targeted-manuscript-validation}

\begin{Shaded}
\begin{Highlighting}[]
\ExtensionTok{pnpm} \AttributeTok{{-}{-}dir}\NormalTok{ open{-}source/aeon/docs/ebooks/145{-}log{-}rolling{-}pipelined{-}prefill/companion{-}tests exec vitest run }\DataTypeTok{\textbackslash{}}
\NormalTok{  src/manuscript{-}flagship{-}structure.test.ts }\DataTypeTok{\textbackslash{}}
\NormalTok{  src/manuscript{-}flagship{-}artifact{-}consistency.test.ts }\DataTypeTok{\textbackslash{}}
\NormalTok{  src/manuscript{-}flagship{-}hardening.test.ts}
\end{Highlighting}
\end{Shaded}

\subsubsection{Supporting scheduler and protocol
checks}\label{supporting-scheduler-and-protocol-checks}

\begin{Shaded}
\begin{Highlighting}[]
\ExtensionTok{pnpm} \AttributeTok{{-}{-}dir}\NormalTok{ open{-}source/aeon/docs/ebooks/145{-}log{-}rolling{-}pipelined{-}prefill/companion{-}tests exec vitest run }\DataTypeTok{\textbackslash{}}
\NormalTok{  src/dag{-}completeness.test.ts }\DataTypeTok{\textbackslash{}}
\NormalTok{  src/pipeline{-}formulas.test.ts }\DataTypeTok{\textbackslash{}}
\NormalTok{  src/pipeline{-}topology.test.ts }\DataTypeTok{\textbackslash{}}
\NormalTok{  src/deficit{-}evidence.test.ts }\DataTypeTok{\textbackslash{}}
\NormalTok{  src/flow{-}protocol.test.ts }\DataTypeTok{\textbackslash{}}
\NormalTok{  src/ch17{-}gate1{-}wallclock{-}figure.test.ts }\DataTypeTok{\textbackslash{}}
\NormalTok{  src/ch17{-}gate2{-}protocol{-}corpus{-}figure.test.ts}
\end{Highlighting}
\end{Shaded}

\subsubsection{Packaged submission
helper}\label{packaged-submission-helper}

\begin{Shaded}
\begin{Highlighting}[]
\FunctionTok{bash}\NormalTok{ open{-}source/aeon/docs/ebooks/145{-}log{-}rolling{-}pipelined{-}prefill/arxiv{-}submission{-}flagship/prepare{-}arxiv.sh}
\end{Highlighting}
\end{Shaded}

The flagship uses the same checked-in companion artifact directory as
the full chapter. It does not define a second artifact format. The
difference is only the manuscript selection, the narrowed claim surface,
and the dedicated build and packaging entrypoints.

\subsection{References}\label{references}

{[}1{]} G. M. Amdahl, ``Validity of the Single Processor Approach to
Achieving Large-Scale Computing Capabilities,'' \emph{AFIPS Spring Joint
Computer Conference}, 30:483-485, 1967.

{[}2{]} J. L. Gustafson, ``Reevaluating Amdahl's Law,''
\emph{Communications of the ACM}, 31(5):532-533, 1988.

{[}3{]} G. Kahn, ``The Semantics of a Simple Language for Parallel
Programming,'' \emph{IFIP Congress}, 1974.

{[}4{]} J. B. Dennis, D. P. Misunas, ``A Preliminary Architecture for a
Basic Data-Flow Processor,'' \emph{Proceedings of the 2nd Annual
Symposium on Computer Architecture}, 1975.

{[}5{]} I. E. Sutherland, ``Micropipelines,'' \emph{Communications of
the ACM}, 32(6):720-738, 1989.

{[}6{]} J. Dean, S. Ghemawat, ``MapReduce: Simplified Data Processing on
Large Clusters,'' \emph{OSDI}, 2004.

{[}7{]} L. Lamport, \emph{Specifying Systems: The TLA+ Language and
Tools for Hardware and Software Engineers}, Addison-Wesley, 2002.

{[}8{]} M. Belshe, R. Peon, M. Thomson, ``Hypertext Transfer Protocol
Version 2 (HTTP/2),'' RFC 7540, 2015.

{[}9{]} J. Iyengar, M. Thomson, ``QUIC: A UDP-Based Multiplexed and
Secure Transport,'' RFC 9000, 2021.

{[}10{]} T. W. Buley, ``Aeon Clockwork: A Unified Probability Engine
with Immanent Self-Verification,'' open-source implementation, 2026.
\url{https://github.com/forkjoin-ai/aeon-clockwork}

\end{document}
